\documentclass{article}


\usepackage{hyperref}
\usepackage{graphicx}
\usepackage{cite}
%\VignetteIndexEntry{Using Animal} 

\hypersetup{
  pdfauthor={Matti Pastell},
  colorlinks=TRUE,
  linkcolor=black,
  citecolor=blue,
  urlcolor=blue
}


%Sivujen otsikot
\usepackage{fancyhdr}
\pagestyle{fancy}
\renewcommand{\sectionmark}[1]{%
\markright{\thesection\ #1}}
\fancyhf{} % delete current header and footer
\fancyhead[LE,RO]{\bfseries\thepage}
\fancyhead[LO]{\bfseries\rightmark}
\fancyhead[RE]{\bfseries\leftmark}
\renewcommand{\headrulewidth}{0.5pt}
\renewcommand{\footrulewidth}{0pt}
\addtolength{\headheight}{0.5pt} % space for the rule
\fancypagestyle{plain}{%
\fancyhead{} % get rid of headers on plain pages
\renewcommand{\headrulewidth}{0pt} % and the line
}

\newcommand{\HRule}{\rule{\linewidth}{0.5mm}}


% \usepackage{graphics}

%Ei sisennyst�
%\setlength\parindent{0in}

% ��kk�set
%\usepackage[utf8]{inputenc}
%\usepackage[T1]{fontenc}
%\usepackage{times}

\usepackage{Sweave}
\begin{document}

%\begin{titlepage}
%\ThisULCornerWallPaper{0.3}{logo}


%\includegraphics[width=0.6\textwidth]{logo}\\[2cm]

\vspace*{1cm}
%\hspace{5cm}


\noindent \textsc{\Large Analyze Animal Behavior Data Using:}\\[0.1cm]
% Title
\HRule \\[0.2cm]
\begin{center}
{ \huge \bfseries Animal}\\[0.4cm]
\HRule \\[5cm]
\end{center}

%\vfill

\noindent \textsc{\Large Matti Pastell}\\[0.5cm]
\noindent \textsc{\Large Department of Agrotechnology \\[0.1cm] \& Research
  Center of Animal Welfare}\\[0.5cm]

%\noindent \url{http://www.mm.helsinki.fi/~mpastell}


%\AddToShipoutPicture{\includegraphics[height=0.5\paperheight]{MMM}}
%\ThisLRCornerWallPaper{0.6}{MMM}

\end{titlepage}


%\begin{thebibliography}{9}
 % \bibitem{}
%\end{thebibliography}


\title{
\noindent \textsc{\Large Analyze Animal Behavior Data Using:}\\
% Title
\HRule \\
\begin{center}
{ \huge \bfseries Animal}\\
\HRule \\[0.2cm]
\end{center}
}

\author{
\Large Matti Pastell\\
\url{matti.pastell@helsinki.fi}\\[0.5cm]
Department of Agrotechnology \& \\
Center of Animal Welfare, \textit{University of Helsinki}
}

  % Matti Pastell \\ matti.pastell@helsinki.fi}

\maketitle




\section{Introduction}
The \texttt{Animal} package is a collection of functions for analyzing animal
behavior data originating from a variety of sources. The package was
originally created to analyze data files from CowLog (open source
software for coding behaviors from digital video), but the
functionality has been extended to cover also other data sources.
This document briefly describes some the key features of the package
via simple examples. \\
Animal has been created in the Research Centre of Animal Welfare
\footnote{\url{http://http://www.vetmed.helsinki.fi/hyvinvointikeskus/english/index.htm}}.


\section{CowLog data}
The package has basic analysis functions for analyzing time coded
behavioral data coded with CowLog\cite{cowlog}
\footnote{\url{http://www.mm.helsinki.fi/~mpastell/CowLog}}.
The main function in the package is \textit{cowAnalyze}. The function
takes the data file name, labels of the codes and the type (event or
state) of the codes as inputs, and gives a summary table and plot of
the results as output. The function also removes double state
(duplicated) errors for state codes.  

Here is a short example on how to use the function in R:
Analyze CowLog datafile named calf1.bh1: Define codes 1-3 and give them names walking, standing and lying. The descriptive statistics will appear on screen and they are also saved to variable analyzed.  
\begin{Schunk}
\begin{Sinput}
> library(Animal)
> analyzed <- cowAnalyze(file = "calf1.bh1", states = c(1, 
+     2, 3), state.names = c("walking", "standing", "lying"))